\documentclass[11pt,a4paper]{article}
\usepackage[margin=25mm]{geometry}
\usepackage[T1]{fontenc}
\usepackage[utf8]{inputenc}
\usepackage{lmodern,microtype}
\usepackage{amsmath,amssymb,booktabs,longtable}
\usepackage{hyperref}
\title{Research Logbook: SmartWeld OSLW Reconstruction and Physics-Informed ML}
\author{Andrey Ni}
\date{30 September 2026}
\begin{document}
\maketitle
\section*{Purpose}
This logbook records source provenance, reconstruction decisions, numerical evidence,
open problems, and research tasks for reconstruction of the SmartWeld 3.0 OSLW model
distributed in the 12 March 2012 MATLAB M-files release.
\section*{Evidence classes}
\begin{description}
\item[Source] Directly supported by recovered SmartWeld source or official distribution.
\item[Regression] Reproduced against an original SmartWeld GUI output.
\item[Reconstruction] Independently implemented from source-derived equations.
\item[Inference] Interpretation requiring explicit qualification.
\item[Proposal] Future work not yet validated.
\end{description}
\section{Research problem}
The project investigates whether a historical physics-based welding model can be
recovered from legacy scientific software, independently reproduced, and exposed as
a reproducible computational component for modern machine-learning and optimization.
The immediate target is OSLW (Optimized Single-pass Laser Welding). The longer-term
objective is a physics-informed workflow for laser processing and, after independent
validation, DED.
\section{Current evidence status}
\begin{longtable}{p{0.27\textwidth}p{0.18\textwidth}p{0.47\textwidth}}
\toprule Task & Status & Evidence\\ \midrule
Official 12 Mar 2012 M-files archive & verified & Official SmartWeld SourceForge distribution; checksum recorded.\\
Standalone executable & verified & Supplied SmartWeld-1March2012.exe.\\
OSLW call chain & recovered & \texttt{app\_specific\_data} $\rightarrow$ \texttt{schedulermodels} $\rightarrow$ \texttt{queryfunc}.\\
Forward equations & reconstructed & Normalization, material constants, gas coefficients, penetration, ETE, melting efficiency, width and sensitivity.\\
Forward validation & verified for 2 cases & Historical GUI outputs reproduced to printed precision.\\
Exact inverse GUI reproduction & incomplete & \texttt{computeScheduleArray} depends on contour generation; \texttt{ccontours.m} remains unresolved.\\
ML model & proposed & No ML performance result is claimed.\\
DED extension & proposed & OSLW is not treated as a validated DED model.\\
\bottomrule
\end{longtable}
\section{Recovered OSLW equations}
The source normalizes process variables:
\begin{equation}
s_P=P/1600,\qquad s_V=V/120,\qquad s_d=d/0.0294.
\end{equation}
For 304 stainless steel:
\begin{equation}
\alpha_s=5.18/110,\qquad \Delta h_s=9.41/11.8.
\end{equation}
For argon:
\begin{align}
C_1&=0.9015503591510303,&C_2&=0.6327730407815897,&C_3&=0.2582188708198456,\\
C_8&=29.32835498219138,&C_9&=0.2903225193306445,&C_{10}&=0.3160507149639552.
\end{align}
Penetration:
\begin{equation}
p_s=\frac{\alpha_s C_8(s_P/[s_d^{C_{10}}s_V^{C_9}])}{5}.
\end{equation}
Energy-transfer efficiency:
\begin{equation}
\eta_{\mathrm{ET},s}=
\frac{C_1-C_2^{\pi/[2\arctan(C_3s_d/p_s)]}}{0.952}.
\end{equation}
The source then evaluates
\begin{equation}
Ry_s=\frac{s_P\eta_{\mathrm{ET},s}s_V}{\alpha_s^2\Delta h_s},
\qquad
Ry=Ry_s\frac{1600\cdot120\cdot0.952}{110^2\cdot11.8},
\end{equation}
and
\begin{equation}
\eta_m=0.48-0.29e^{-Ry/6.82}-0.17e^{-Ry/58.8}.
\end{equation}
Area is evaluated as
\begin{equation}
A_s=\frac{Ry\eta_m\alpha_s^2}{s_V^2},
\qquad A=A_s(110/120)^2.
\end{equation}
\section{Regression cases}
\begin{center}
\begin{tabular}{lcc}
\toprule Quantity & Case 1 & Case 2\\ \midrule
Material & 304 stainless & 304 stainless\\
Gas & Ar & Ar\\
Power (W) & 685.000000 & 786.666667\\
Speed (mm/s) & 53.082625 & 52.839095\\
Spot (cm) & 0.011800 & 0.011800\\
Target penetration (mm) & 1.000 & 1.150\\
Reported penetration (mm) & 0.999860 & 1.149791\\
Width (mm) & 0.625687 & 0.678073\\
ETE & 0.679181 & 0.718454\\
Melting efficiency & 0.447787 & 0.457256\\
Sensitivity ($\mu$m/W) & 1.459649 & 1.461599\\
\bottomrule
\end{tabular}
\end{center}
The independent forward reconstruction reproduces these values to the precision
printed by the historical GUI outputs.
\section{Inverse-solver issue}
The recovered GUI path calls \texttt{computeScheduleArray}. In depth-only mode it
selects the midpoint of a returned process contour before evaluating
\texttt{queryfunc}. Therefore the historical inverse result cannot be replaced by
an assumed continuous optimization objective. Recovering \texttt{ccontours.m} is
the next critical source-recovery target.
\section{Research tasks}
\begin{enumerate}
\item Source completeness: recover \texttt{ccontours.m}, verify inventory, freeze archive hash.
\item Numerical reproducibility: deterministic regression tests and broader parameter grid.
\item Physics-based dataset: source-valid DOE with complete metadata.
\item Machine learning: baseline and physics-informed surrogates with uncertainty estimation.
\item Experimental transfer: map to laser/DED hardware and validate experimentally.
\end{enumerate}
\section{Publication decision rules}
The first preprint may claim source-level reconstruction of the OSLW forward model,
two-case numerical reproduction, documented inverse-workflow analysis, and a
reproducible framework for later physics-informed ML. It should not claim complete
inverse-GUI reproduction before \texttt{ccontours.m} is recovered, experimentally
validated ML improvement before ML experiments exist, or a validated DED model from
OSLW alone.
\section{Change log}
\begin{longtable}{p{0.15\textwidth}p{0.12\textwidth}p{0.63\textwidth}}
\toprule Date & Version & Entry\\ \midrule
2026-09-30 & 0.1 & Initial publication logbook; forward reconstruction and two-case regression established; inverse contour reconstruction remains open.\\
\bottomrule
\end{longtable}
\end{document}